\chapter*{2. Evidencia y literatura relevante}

\section*{2.1. Evidencia}

La principal hipótesis que se revisa en esta investigación es que el cambio climático --y, en particular, el estrés térmico-- tiene un efecto importante sobre el mercado internacional del café, en tanto se trata de un bien agrícola que requiere de condiciones climáticas moderadas y controladas para su desarrollo.

En tal sentido, la [Figura N°1] muestra la evolución mensual de la temperatura promedio, desde 1965 hasta 2019, para todo el mundo y para las zonas específicas dedicadas a la producción de café. Se observa [DESARROLLAR].

De esta manera, sí existe un diferencial entre la temperatura en los países cafetaleros, que es más elevada por estás en la franja ecuatorial, y el resto del mundo. Por otro lado, tal y como muestra la [Figura N°2], esta diferencia no es significativa entre grupos de países exportadores.

Ahora bien, para efectos de medir el impacto del estrés térmico sobre el mercado, la variable relevante para el análisis es la acumulación de temperaturas extremas. La construcción de esta variable es explicada con detalle en la sección [COMPLETAR], y la [Figura N°3] presenta su evolución en el tiempo. Nuevamente, se observa una diferencia importante según si se trata de un país cafetalero o no. En promedio, en el mismo mes, un país exportador de café puede acumular hasta [22] días de calor, alejándose notoriamente del promedio mundial de 15 días.

Por otro lado, si bien las exportaciones de café aumentan en cada año, se observa una diferencia importante entre la evolución de países líderes y seguidores. En resumen, mientras que en la segunda mitad del siglo XIX los países líderes concentraban el XX\% del mercado, hoy alcanzan el XX\%. La [Figura N°4] refleja esta importante diferencia en la tendencia de cada grupo exportador.

Por último, y volviendo a la hipótesis principal, se observa una correlación negativa entre el aumento del estrés térmico y las exportaciones globales (Figura N°4). Además, esta correlación es más pronunciada para los países seguidores (Figura N°5). Lo mismo se observa para el precio del café (Figura N°6).

\section*{2.2. Literatura relevante}

La literatura económica relativa al impacto del cambio climático ha evolucionado significativamente en los últimos años, transitando desde enfoques basados en la correlación de variables geográficas fijas hacia metodologías dinámicas que explotan variaciones temporales de fenómenos climáticos [CITAR]. El método convencional por mucho tiempo se ha basado en la estimación con datos de panel, utilizando fluctuaciones interanuales de la temperatura a nivel local, conocidas como choques de clima [CITAR]. Este enfoque se centra en choques climáticos de corta duración; con frecuencia, se utilizan efectos fijos temporales para eliminar impactos comunes agregados [CITAR].

Sin embargo, recientemente se ha argumentado que un enfoque más sólido para capturar los efectos persistentes del cambio climático, así como las respuestas de adaptación, reside en metodologías de diferencias a largo plazo (\textit{long differences}), que comparan cambios en las tendencias climáticas a lo largo de períodos más extensos (opr ejemplo, décadas) [CITAR]. Esta comparación de largo plazo permite evaluar si la adaptación realizada por los agentes económicos mitiga el daño observado en el corto plazo [CITAR]. Por ejemplo, algunas investigaciones sugieren que, si bien la adaptación es un factor que mitiga los efectos negativos, esta capacidad puede ser limitada o nula para contrarrestar los efectos negativos del calor extremo en la productividad de bienes agrícolas [CITAR]. Esto se debería a costos de ajuste importantes y a la dependecia de cada cultivo a condiciones geográficas específicas.

Un enfoque estructural para este problema ofrece ventajas importantes en el análisis económico, pues racionaliza la interacción estratégica de los agentes, lo cual es fundamental en mercados oligopolísticos como el del café. Los modelos estructurales permiten descomponer el precio en costo marginal y \textit{markup}, e identificar el grado de interdependencia en el mercado [CITAR]. En este caso, permite ir más allá de la correlación de variables específicas y evaluar cómo los cambios en las condiciones de mercado --en este caso, choques climáticos-- reconfiguran el equilibrio estrat´gico y el poder de mercado, lo cual es esencial para el análisis de bienestar [CITAR].

Esta investigación utiliza una estimación estructural para cuantificar el efecto del cambio climático sobre el mercado del café, específicamente a través de un modelo de competencia secuencial en cantidades, estableciendo jerarquías entre los países exportadores, de cara a la comercialización de un bien homogéneo.

La aproximación estructural al problema de estimación de impacto del cambio climático permite simular escenarios contrafactuales, a fin de evaluar el impacto en el bienestar social y la distribución entre oferentes y demandantes.

Adicionalmente, se debe tener en consideración que este no es el primer trabajo que estima el mercado internacional del café de forma estructural. Karp y Perloff (1993) desarrollaron un modelo de oligopolio para estudiar si los dos mayores exportadores de café de la época, Brazil y Colombia, operaban como tomadores de precios, oligopolistas, o como partes de un equilibrio de colusión tácita. Sus resultados indicaron un comportamiento relativamente competitivo, similar a lo encontrado por Bettendorf y Verboren (2000), cuyo objetivo de análisis fue estimar el traspaso del precio global al precio de retail en Países Bajos; y por Gilbert (2007), quien encontró un aumento de la intensidad competitiva en el mercado de comercio internacional del café y en el mercado de retail del mismo bien.

Otros trabajos han utilizado estimaciones estructurales del café y otros bienes agrícolas para cuantificar el efecto de choques exógenos. Este es el caso de Roberts y Schlenker (2010), quienes utilizaron desvíos en la tendencia del rendimiento de los cultivos de bienes agrícolas -bajo el supuesto de que estos desvíos son explicados por choques climáticos exógenos- para medir el impacto del Mandato de Etanol de Estados Unidos, que exigía que alrededor del 5\% de la producción mundial de calorías de maíz, trigo, arroz y soja se destine a la producción de etanol. Como resultado, estimaron que los precios mundiales de los alimentos aumentarían en torno a un 30\% y que el excedente mundial del consumo de alimentos disminuiría en 155 mil millones de dólares anuales.

En esa misma línea, el trabajo de Igami (2015) constituye otro importante ejemplo, cuyo modelo oligopólico es utilizado para cuantificar el perjuicio generado por la cartelización de los países miembros de la ICO entre 1965 y 1989 sobre el bienestar social del mercado.

Por último, los trabajos de Mehta y Chavas (2005), Maurice y Davis (2011), Covindassamy, Robe y Wallen (2017), y Lanna, Mattos y Sylvia (2017), entre otros, estudian la incertidumbre en las series temporales de los precios de bienes agrícolas como el café, a fin de explicar las razones detrás de su volatilidad.

En este contexto, la presente investigación toma una herramienta ya consolidada en el estudio de este \textit{soft commodity}, a fin de estimar el mercado internacional del café a través de un modelo de competencia en cantidades, y cuantificar el impacto del cambio climático, así como la distribución de la pérdida de bienestar social asociada. Al mismo tiempo, este trabajo propone una estructura de competencia secuencial entre países líderes y seguidores, en contraste con los modelos de competencia simultánea hasta ahora desarrollados. Este enfoque permite capturar de mejor manera la racionalidad detrás de los agentes que interactúan en el mercado, pues distingue entre grandes actores, caracterizados por un nivel de comercialización notablemente mayor y una trayectoria competitiva histórica, y agentes seguidores, cuya presencia en el mercado es relevante, mas no determinante para efectos de incidir sobre el equilibrio competitivo. 

